\documentclass[11pt,a4paper]{article}
\usepackage[margin=0.85in]{geometry}
\usepackage{amsmath,amssymb,booktabs,array,graphicx,float,hyperref}
\usepackage{caption}
\usepackage{microtype}
\hypersetup{colorlinks=true,linkcolor=blue!55!black,urlcolor=blue!55!black,citecolor=blue!55!black}
\setlength{\parskip}{0.55em}
\setlength{\parindent}{0pt}
\title{\textbf{Zhang-Style Forecast Evaluation for Dow 30, S\&P 100, and S\&P 500 GNNHAR Runs}}
\author{Generated from local Volatility outputs}
\date{June 18, 2026}
\begin{document}
\maketitle

\begin{abstract}
This report summarizes the post-run evidence generated from the saved Dow 30, S\&P 100,
and S\&P 500 GNNHAR forecasts.  The objective is not to retrain models, but to read the
local prediction arrays and compute the statistics used in Zhang et al. (2024a)'s GNNHAR paper:
HAR-normalized MSE and QLIKE loss ratios, model confidence sets (Hansen et al., 2011), calm-versus-turbulent
regime tables, pairwise QLIKE Diebold-Mariano comparisons (Diebold and Mariano, 1995), FVU nonlinearity diagnostics,
forecast error and forecast ratio boxplot sources, graph summaries, and smoothing diagnostics.
The S\&P 500 analysis is now tied to the actual AutoDL run source rather than the smaller
Google Drive upload panel.
\end{abstract}


\section{Notation}
Let \(v_{i,t}\) denote the realized variance or realized volatility target for ticker \(i\) on
date \(t\), and let \(\widehat v_{i,t}^{(m)}\) be the forecast from model \(m\).
The benchmark model is \(HAR_{\mathrm{M}}\), the HAR model estimated under MSE.  The loss ratios are
\[
R^{MSE}_m =
\frac{\frac{1}{NT}\sum_{i,t}(v_{i,t}-\widehat v_{i,t}^{(m)})^2}
{\frac{1}{NT}\sum_{i,t}(v_{i,t}-\widehat v_{i,t}^{(HAR_{\mathrm{M}})})^2},
\quad
R^{QL}_m =
\frac{\frac{1}{NT}\sum_{i,t}QL(v_{i,t},\widehat v_{i,t}^{(m)})}
{\frac{1}{NT}\sum_{i,t}QL(v_{i,t},\widehat v_{i,t}^{(HAR_{\mathrm{M}})})}.
\]
The QLIKE loss is
\[
QL(y,\widehat y)=\frac{y}{\widehat y}-\log\!\left(\frac{y}{\widehat y}\right)-1.
\]
A ratio below one means lower average forecast loss than \(HAR_{\mathrm{M}}\).  \(GHARkH\) denotes a
linear graph-HAR model using \(k\)-hop neighbors; \(GNNHARkL\) denotes a \(k\)-layer graph
neural HAR model.  The suffix \(_IV\) indicates that lagged implied volatility is included.
MCS denotes the Hansen--Lunde--Nason model confidence set at the \(5\%\) level.  DM denotes
the QLIKE Diebold--Mariano comparison; in the generated files, a positive statistic favors the
candidate model listed in the comparison.


\section{Data Source Alignment and the 397-versus-449 Issue}

The earlier discrepancy came from comparing the S\&P 500 AutoDL forecasts against the wrong raw panel.  The formal AutoDL run used the local equivalent of \texttt{data/scale\_experiment/sp500}, whose RV, IV, and return panels each contain 503 tickers before coverage filtering.  The model run selected 449 tickers.  The smaller 397-column RV/IV panel belongs to the Google Drive upload copy and is not the correct source for the formal S\&P 500 run.  After rebinding the analysis script to the AutoDL metadata source, all 449 model tickers match the S\&P 500 RV, IV, and return panels.

\begin{table}[H]
\centering
\caption{Run-level sample summary}
\label{tab:sample-summary}
\scriptsize
\resizebox{\textwidth}{!}{%
\begin{tabular}{lrrrrrrr}
\toprule
Universe & Test dates & Tickers & Models & Test start & Test end & Raw source rule & RV missing \\
\midrule
dow30 & 234 & 30 & 8 & 2025-07-07 & 2026-06-09 & aligned\_wide\_multihop\_rerun & 0 \\
sp100 & 234 & 91 & 32 & 2025-07-07 & 2026-06-09 & run\_config\_data\_paths\_data\_dir & 0 \\
sp500 & 223 & 449 & 36 & 2025-07-14 & 2026-06-01 & autodl\_run\_metadata\_args\_data\_dir & 0 \\
\bottomrule
\end{tabular}
}

\vspace{2pt}\begin{minipage}{0.98\textwidth}\footnotesize The table uses the current aligned sample where available. Dow 30 now refers to the 2025-07-07--2026-06-09 aligned wide-window multi-hop rerun; the older 232-date Dow 30 full-model run is archived evidence and should not be described as the current out-of-sample date range. The S\&P 500 source rule is AutoDL run metadata.\end{minipage}
\end{table}

\begin{table}[H]
\centering
\caption{Ticker alignment audit}
\label{tab:ticker-alignment}
\scriptsize
\resizebox{\textwidth}{!}{%
\begin{tabular}{lrrrrrrr}
\toprule
Universe & Panel & Source rule & Panel tickers & Model tickers & Matched & Missing & Extra \\
\midrule
dow30 & local\_raw\_rv\_panel & run\_config\_data\_paths\_data\_dir & 30 & 30 & 30 & 0 & 0 \\
dow30 & local\_raw\_return\_panel & run\_config\_data\_paths\_data\_dir & 30 & 30 & 30 & 0 & 0 \\
dow30 & local\_raw\_iv\_panel & run\_config\_data\_paths\_data\_dir & 30 & 30 & 30 & 0 & 0 \\
sp100 & local\_raw\_rv\_panel & run\_config\_data\_paths\_data\_dir & 91 & 91 & 91 & 0 & 0 \\
sp100 & local\_raw\_return\_panel & run\_config\_data\_paths\_data\_dir & 101 & 91 & 91 & 0 & 10 \\
sp100 & local\_raw\_iv\_panel & run\_config\_data\_paths\_data\_dir & 91 & 91 & 91 & 0 & 0 \\
sp500 & local\_raw\_rv\_panel & autodl\_run\_metadata\_args\_data\_dir & 503 & 449 & 449 & 0 & 54 \\
sp500 & local\_raw\_return\_panel & autodl\_run\_metadata\_args\_data\_dir & 503 & 449 & 449 & 0 & 54 \\
sp500 & local\_raw\_iv\_panel & autodl\_run\_metadata\_args\_data\_dir & 503 & 449 & 449 & 0 & 54 \\
\bottomrule
\end{tabular}
}

\vspace{2pt}\begin{minipage}{0.98\textwidth}\footnotesize Extra tickers are available in the raw panel but excluded by the modeling coverage filter; they are not missing model inputs.\end{minipage}
\end{table}

\section{Implemented Zhang-Style Statistics}

\begin{table}[H]
\centering
\caption{Coverage of Zhang-style statistics}
\label{tab:coverage}
\scriptsize
\resizebox{\textwidth}{!}{%
\begin{tabular}{lrr}
\toprule
Component & Paper reference & Status \\
\midrule
forecast\_loss\_ratios & Table 1 and Table 5 & implemented\_from\_saved\_forecasts \\
weekly\_monthly\_horizons & Table 1, Table 2, Table 3, Table 5 & scope\_gap\_current\_saved\_forecasts\_are\_one\_day\_only \\
model\_confidence\_set & Table 1, Table 2, Table 4, Table 5 & implemented\_from\_per\_ticker\_losses \\
market\_regime\_losses & Table 2 & implemented\_from\_saved\_forecasts \\
alternative\_validation\_size & Table 4 & scope\_gap\_requires\_new\_rolling\_runs \\
forecast\_error\_ratio\_boxplots & Figure 5 & implemented\_as\_source\_tables\_and\_pngs \\
fvu\_nonlinearity & Table 3 and equation (11) & implemented\_with\_regime\_split \\
multi\_hop\_dm & Figure 6 and Appendix E.1 & implemented\_when\_models\_exist \\
mad\_smoothing & Figure 7 and equation (12) & implemented\_when\_graph\_and\_hidden\_or\_forecast\_representations\_exist \\
data\_summary\_and\_correlation & Appendix A and Figure 1 & implemented\_for\_available\_local\_panels \\
graph\_structure & Appendix A.2 & implemented\_when\_graph\_matrices\_exist \\
\bottomrule
\end{tabular}
}

\vspace{2pt}\begin{minipage}{0.98\textwidth}\footnotesize The current saved forecasts are one-day horizon forecasts. Zhang's one-week and one-month tables require separate horizon-specific target construction and reruns. Exact hidden-state MAD requires saved hidden representations; current diagnostics include forecast-level smoothing proxies.\end{minipage}
\end{table}

\section{Main Loss-Ratio Evidence}

Tables below report the complete model set available for each universe in the current one-day-horizon saved forecasts.  These tables correspond to Zhang's Table 1 and Table 5 style, with the added IV variants and the S\&P 500 extension.

\begin{table}[H]
\centering
\caption{Complete loss ratios for DOW30}
\label{tab:dow30-loss}
\tiny
\resizebox{\textwidth}{!}{%
\begin{tabular}{lrrrrrrrr}
\toprule
Model & Family & Depth/hop & IV & Est. & MSE ratio & QL ratio & MCS MSE & MCS QL \\
\midrule
GNNHAR2L\_Q\_IV & GNNHAR & L2 & yes & QLIKE & 0.892 & 0.889 & yes & yes \\
GHAR\_Q\_IV & GHAR & H1 & yes & QLIKE & 0.908 & 0.909 & no & no \\
GNNHAR3L\_Q\_IV & GNNHAR & L3 & yes & QLIKE & 0.916 & 0.912 & no & no \\
GNNHAR1L\_Q\_IV & GNNHAR & L1 & yes & QLIKE & 0.924 & 0.914 & no & no \\
GNNHAR2L\_M\_IV & GNNHAR & L2 & yes & MSE & 0.948 & 0.950 & no & no \\
HAR\_Q\_IV & HAR &  & yes & QLIKE & 0.998 & 0.957 & no & no \\
GNNHAR2L\_Q & GNNHAR & L2 & no & QLIKE & 0.967 & 0.957 & no & no \\
GNNHAR3L\_Q & GNNHAR & L3 & no & QLIKE & 0.997 & 0.971 & no & no \\
GNNHAR1L\_Q & GNNHAR & L1 & no & QLIKE & 0.986 & 0.980 & no & no \\
HAR\_M & HAR &  & no & MSE & 1.000 & 1.000 & no & no \\
GNNHAR1L\_M & GNNHAR & L1 & no & MSE & 0.945 & 1.007 & no & no \\
HAR\_Q & HAR &  & no & QLIKE & 1.090 & 1.013 & no & no \\
GHAR\_Q & GHAR & H1 & no & QLIKE & 0.959 & 1.021 & no & no \\
HAR\_M\_IV & HAR &  & yes & MSE & 0.920 & 1.030 & no & no \\
GNNHAR3L\_M & GNNHAR & L3 & no & MSE & 0.982 & 1.058 & no & no \\
GNNHAR1L\_M\_IV & GNNHAR & L1 & yes & MSE & 0.935 & 1.076 & no & no \\
GHAR\_M\_IV & GHAR & H1 & yes & MSE & 0.868 & 1.137 & yes & no \\
GHAR\_M & GHAR & H1 & no & MSE & 0.927 & 1.140 & no & no \\
GNNHAR2L\_M & GNNHAR & L2 & no & MSE & 0.983 & 1.222 & no & no \\
GNNHAR3L\_M\_IV & GNNHAR & L3 & yes & MSE & 0.995 & 1.233 & no & no \\
\bottomrule
\end{tabular}
}

\vspace{2pt}\begin{minipage}{0.98\textwidth}\footnotesize Ratios are relative to \(HAR_{\mathrm{M}}\). MCS columns indicate inclusion in the model confidence set based on per-ticker losses.\end{minipage}
\end{table}

\begin{table}[H]
\centering
\caption{Complete loss ratios for SP100}
\label{tab:sp100-loss}
\tiny
\resizebox{\textwidth}{!}{%
\begin{tabular}{lrrrrrrrr}
\toprule
Model & Family & Depth/hop & IV & Est. & MSE ratio & QL ratio & MCS MSE & MCS QL \\
\midrule
HAR\_Q\_IV & HAR &  & yes & QLIKE & 0.981 & 0.933 & yes & yes \\
GHAR\_Q\_IV & GHAR & H1 & yes & QLIKE & 0.981 & 0.933 & yes & yes \\
GNNHAR2L\_Q\_IV & GNNHAR & L2 & yes & QLIKE & 0.982 & 0.933 & no & yes \\
GNNHAR5L\_Q\_IV & GNNHAR & L5 & yes & QLIKE & 0.981 & 0.934 & yes & yes \\
GNNHAR1L\_Q\_IV & GNNHAR & L1 & yes & QLIKE & 0.981 & 0.935 & no & yes \\
GNNHAR4L\_Q\_IV & GNNHAR & L4 & yes & QLIKE & 0.986 & 0.936 & no & yes \\
GNNHAR3L\_Q\_IV & GNNHAR & L3 & yes & QLIKE & 0.981 & 0.940 & no & no \\
GHAR\_Q & GHAR & H1 & no & QLIKE & 1.003 & 0.976 & no & no \\
GNNHAR5L\_Q & GNNHAR & L5 & no & QLIKE & 1.003 & 0.977 & no & no \\
GHAR2H\_M\_IV & GHAR & H2 & yes & MSE & 0.965 & 0.977 & yes & no \\
GNNHAR3L\_Q & GNNHAR & L3 & no & QLIKE & 1.003 & 0.977 & no & no \\
GHAR\_M\_IV & GHAR & H1 & yes & MSE & 0.965 & 0.978 & yes & no \\
GHAR3H\_M\_IV & GHAR & H3 & yes & MSE & 0.965 & 0.978 & yes & no \\
GNNHAR4L\_Q & GNNHAR & L4 & no & QLIKE & 1.003 & 0.978 & no & no \\
HAR\_M\_IV & HAR &  & yes & MSE & 0.965 & 0.979 & yes & no \\
GNNHAR2L\_Q & GNNHAR & L2 & no & QLIKE & 1.004 & 0.979 & no & no \\
HAR\_Q & HAR &  & no & QLIKE & 1.006 & 0.980 & no & no \\
GNNHAR1L\_Q & GNNHAR & L1 & no & QLIKE & 1.004 & 0.980 & no & no \\
GNNHAR1L\_M\_IV & GNNHAR & L1 & yes & MSE & 0.972 & 0.987 & yes & no \\
GNNHAR3L\_M\_IV & GNNHAR & L3 & yes & MSE & 0.968 & 0.990 & yes & no \\
GHAR\_M & GHAR & H1 & no & MSE & 0.997 & 0.993 & no & no \\
GNNHAR5L\_M\_IV & GNNHAR & L5 & yes & MSE & 0.980 & 0.994 & no & no \\
GNNHAR4L\_M\_IV & GNNHAR & L4 & yes & MSE & 0.979 & 0.996 & no & no \\
GNNHAR1L\_M & GNNHAR & L1 & no & MSE & 1.000 & 0.997 & no & no \\
GHAR2H\_M & GHAR & H2 & no & MSE & 0.998 & 0.998 & no & no \\
GNNHAR2L\_M & GNNHAR & L2 & no & MSE & 1.000 & 0.999 & no & no \\
GNNHAR3L\_M & GNNHAR & L3 & no & MSE & 1.001 & 1.000 & no & no \\
HAR\_M & HAR &  & no & MSE & 1.000 & 1.000 & no & no \\
GHAR3H\_M & GHAR & H3 & no & MSE & 0.998 & 1.001 & no & no \\
GNNHAR4L\_M & GNNHAR & L4 & no & MSE & 1.004 & 1.003 & no & no \\
GNNHAR5L\_M & GNNHAR & L5 & no & MSE & 1.003 & 1.010 & no & no \\
GNNHAR2L\_M\_IV & GNNHAR & L2 & yes & MSE & 0.989 & 1.012 & no & no \\
\bottomrule
\end{tabular}
}

\vspace{2pt}\begin{minipage}{0.98\textwidth}\footnotesize Ratios are relative to \(HAR_{\mathrm{M}}\). MCS columns indicate inclusion in the model confidence set based on per-ticker losses.\end{minipage}
\end{table}

\begin{table}[H]
\centering
\caption{Complete loss ratios for SP500}
\label{tab:sp500-loss}
\tiny
\resizebox{\textwidth}{!}{%
\begin{tabular}{lrrrrrrrr}
\toprule
Model & Family & Depth/hop & IV & Est. & MSE ratio & QL ratio & MCS MSE & MCS QL \\
\midrule
GHAR\_M\_IV & GHAR & H1 & yes & MSE & 0.966 & 0.973 & yes & yes \\
HAR\_M\_IV & HAR &  & yes & MSE & 0.966 & 0.974 & yes & no \\
GHAR2H\_M\_IV & GHAR & H2 & yes & MSE & 0.986 & 0.991 & no & no \\
GHAR\_M & GHAR & H1 & no & MSE & 1.000 & 0.998 & no & no \\
GHAR2H\_M & GHAR & H2 & no & MSE & 1.001 & 0.998 & no & no \\
GHAR3H\_M & GHAR & H3 & no & MSE & 1.001 & 0.998 & no & no \\
HAR\_M & HAR &  & no & MSE & 1.000 & 1.000 & no & no \\
GHAR3H\_M\_IV & GHAR & H3 & yes & MSE & 1.002 & 1.020 & no & no \\
GHAR3H\_Q\_IV & GHAR & H3 & yes & QLIKE & 88.253 & 2.621 & no & no \\
GHAR2H\_Q\_IV & GHAR & H2 & yes & QLIKE & 88.340 & 2.624 & no & no \\
HAR\_Q\_IV & HAR &  & yes & QLIKE & 124.370 & 2.677 & no & no \\
GHAR\_Q\_IV & GHAR & H1 & yes & QLIKE & 83.373 & 2.698 & no & no \\
GHAR\_Q & GHAR & H1 & no & QLIKE & 161.079 & 2.851 & no & no \\
GHAR3H\_Q & GHAR & H3 & no & QLIKE & 142.460 & 2.895 & no & no \\
GHAR2H\_Q & GHAR & H2 & no & QLIKE & 142.009 & 2.899 & no & no \\
HAR\_Q & HAR &  & no & QLIKE & 180.880 & 2.910 & no & no \\
GNNHAR2L\_Q\_IV & GNNHAR & L2 & yes & QLIKE & 12.851 & 7.802 & no & no \\
GNNHAR4L\_Q\_IV & GNNHAR & L4 & yes & QLIKE & 13.632 & 7.808 & no & no \\
GNNHAR1L\_Q\_IV & GNNHAR & L1 & yes & QLIKE & 13.818 & 7.823 & no & no \\
GNNHAR3L\_Q\_IV & GNNHAR & L3 & yes & QLIKE & 12.930 & 7.824 & no & no \\
GNNHAR5L\_Q\_IV & GNNHAR & L5 & yes & QLIKE & 13.241 & 7.826 & no & no \\
GNNHAR3L\_Q & GNNHAR & L3 & no & QLIKE & 17.289 & 7.854 & no & no \\
GNNHAR5L\_Q & GNNHAR & L5 & no & QLIKE & 17.652 & 7.856 & no & no \\
GNNHAR1L\_Q & GNNHAR & L1 & no & QLIKE & 17.051 & 7.857 & no & no \\
GNNHAR2L\_Q & GNNHAR & L2 & no & QLIKE & 16.106 & 8.152 & no & no \\
GNNHAR4L\_Q & GNNHAR & L4 & no & QLIKE & 19.880 & 9.965 & no & no \\
GNNHAR1L\_M\_IV & GNNHAR & L1 & yes & MSE & 6.261 & 10.912 & no & no \\
GNNHAR5L\_M & GNNHAR & L5 & no & MSE & 6.272 & 10.919 & no & no \\
GNNHAR1L\_M & GNNHAR & L1 & no & MSE & 6.282 & 10.922 & no & no \\
GNNHAR3L\_M\_IV & GNNHAR & L3 & yes & MSE & 6.321 & 10.935 & no & no \\
GNNHAR5L\_M\_IV & GNNHAR & L5 & yes & MSE & 6.361 & 10.951 & no & no \\
GNNHAR2L\_M\_IV & GNNHAR & L2 & yes & MSE & 6.380 & 10.956 & no & no \\
GNNHAR3L\_M & GNNHAR & L3 & no & MSE & 6.466 & 10.987 & no & no \\
GNNHAR4L\_M\_IV & GNNHAR & L4 & yes & MSE & 6.578 & 11.048 & no & no \\
GNNHAR4L\_M & GNNHAR & L4 & no & MSE & 6.743 & 11.098 & no & no \\
GNNHAR2L\_M & GNNHAR & L2 & no & MSE & 7.040 & 11.258 & no & no \\
\bottomrule
\end{tabular}
}

\vspace{2pt}\begin{minipage}{0.98\textwidth}\footnotesize Ratios are relative to \(HAR_{\mathrm{M}}\). MCS columns indicate inclusion in the model confidence set based on per-ticker losses.\end{minipage}
\end{table}

\begin{table}[H]
\centering
\caption{Best models by universe and evaluation loss}
\label{tab:best-models}
\scriptsize
\resizebox{\textwidth}{!}{%
\begin{tabular}{lrrrrrr}
\toprule
Universe & Selection metric & Best model & Ratio & Gain & Family & IV \\
\midrule
dow30 & mse\_ratio\_vs\_HAR\_M & GHAR\_M\_IV & 0.868 & 13.2\% & GHAR & yes \\
dow30 & qlike\_ratio\_vs\_HAR\_M & GNNHAR2L\_Q\_IV & 0.889 & 11.1\% & GNNHAR & yes \\
sp100 & mse\_ratio\_vs\_HAR\_M & HAR\_M\_IV & 0.965 & 3.5\% & HAR & yes \\
sp100 & qlike\_ratio\_vs\_HAR\_M & HAR\_Q\_IV & 0.933 & 6.7\% & HAR & yes \\
sp500 & mse\_ratio\_vs\_HAR\_M & HAR\_M\_IV & 0.966 & 3.4\% & HAR & yes \\
sp500 & qlike\_ratio\_vs\_HAR\_M & GHAR\_M\_IV & 0.973 & 2.7\% & GHAR & yes \\
\bottomrule
\end{tabular}
}

\vspace{2pt}\begin{minipage}{0.98\textwidth}\footnotesize The best QLIKE model differs by universe: Dow 30 favors a nonlinear IV model, S\&P 100 favors HAR with IV, and S\&P 500 favors GHAR with IV.\end{minipage}
\end{table}

\section{MCS and DM Interpretation}

The MCS results are stricter than simple loss ranking.  Dow 30 has a clear QLIKE winner \(GNNHAR2L\_Q\_IV\).  S\&P 100 has a broad IV-trained QLIKE confidence set, indicating that several IV variants are statistically hard to separate.  S\&P 500 is much more linear in the current run: \texttt{GHAR\_M\_IV} is the QLIKE MCS winner, while most QLIKE-trained neural specifications have very poor loss ratios.

\begin{table}[H]
\centering
\caption{QLIKE DM comparison summary for DOW30}
\label{tab:dow30-dm}
\tiny
\resizebox{\textwidth}{!}{%
\begin{tabular}{lrrrrrr}
\toprule
Type & Base & Candidate & Cand./base & Gain & p<0.05 share & Positive share \\
\midrule
graph\_linear & HAR\_M & GHAR\_M & 1.140 & -14.0\% & 0.0\% & 16.7\% \\
graph\_linear & HAR\_M\_IV & GHAR\_M\_IV & 1.104 & -10.4\% & 3.3\% & 30.0\% \\
graph\_linear & HAR\_Q & GHAR\_Q & 1.008 & -0.8\% & 6.7\% & 63.3\% \\
graph\_linear & HAR\_Q\_IV & GHAR\_Q\_IV & 0.950 & 5.0\% & 20.0\% & 83.3\% \\
nonlinear\_one\_hop & GHAR\_M & GNNHAR1L\_M & 0.883 & 11.7\% & 10.0\% & 76.7\% \\
nonlinear\_one\_hop & GHAR\_M\_IV & GNNHAR1L\_M\_IV & 0.946 & 5.4\% & 0.0\% & 66.7\% \\
nonlinear\_one\_hop & GHAR\_Q & GNNHAR1L\_Q & 0.960 & 4.0\% & 0.0\% & 56.7\% \\
nonlinear\_one\_hop & GHAR\_Q\_IV & GNNHAR1L\_Q\_IV & 1.005 & -0.5\% & 10.0\% & 40.0\% \\
gnn\_depth & GNNHAR1L\_M & GNNHAR2L\_M & 1.214 & -21.4\% & 56.7\% & 10.0\% \\
gnn\_depth & GNNHAR1L\_M\_IV & GNNHAR2L\_M\_IV & 0.883 & 11.7\% & 43.3\% & 90.0\% \\
gnn\_depth & GNNHAR1L\_Q & GNNHAR2L\_Q & 0.976 & 2.4\% & 13.3\% & 73.3\% \\
gnn\_depth & GNNHAR1L\_Q\_IV & GNNHAR2L\_Q\_IV & 0.973 & 2.7\% & 23.3\% & 80.0\% \\
gnn\_depth & GNNHAR2L\_M & GNNHAR3L\_M & 0.866 & 13.4\% & 13.3\% & 83.3\% \\
gnn\_depth & GNNHAR2L\_M\_IV & GNNHAR3L\_M\_IV & 1.297 & -29.7\% & 66.7\% & 3.3\% \\
gnn\_depth & GNNHAR2L\_Q & GNNHAR3L\_Q & 1.015 & -1.5\% & 10.0\% & 26.7\% \\
gnn\_depth & GNNHAR2L\_Q\_IV & GNNHAR3L\_Q\_IV & 1.026 & -2.6\% & 3.3\% & 30.0\% \\
\bottomrule
\end{tabular}
}

\vspace{2pt}\begin{minipage}{0.98\textwidth}\footnotesize A candidate/base ratio below one means the candidate has lower average QLIKE than the base.  Positive DM share is the share of tickers where the DM statistic favors the candidate.\end{minipage}
\end{table}

\begin{table}[H]
\centering
\caption{QLIKE DM comparison summary for SP100}
\label{tab:sp100-dm}
\tiny
\resizebox{\textwidth}{!}{%
\begin{tabular}{lrrrrrr}
\toprule
Type & Base & Candidate & Cand./base & Gain & p<0.05 share & Positive share \\
\midrule
graph\_linear & HAR\_M & GHAR\_M & 0.993 & 0.7\% & 22.0\% & 47.3\% \\
graph\_linear & HAR\_M\_IV & GHAR\_M\_IV & 0.999 & 0.1\% & 12.1\% & 57.1\% \\
graph\_linear & HAR\_Q & GHAR\_Q & 0.996 & 0.4\% & 6.6\% & 56.0\% \\
graph\_linear & HAR\_Q\_IV & GHAR\_Q\_IV & 1.001 & -0.1\% & 4.4\% & 50.5\% \\
nonlinear\_one\_hop & GHAR\_M & GNNHAR1L\_M & 1.004 & -0.4\% & 15.4\% & 40.7\% \\
nonlinear\_one\_hop & GHAR\_M\_IV & GNNHAR1L\_M\_IV & 1.009 & -0.9\% & 18.7\% & 44.0\% \\
nonlinear\_one\_hop & GHAR\_Q & GNNHAR1L\_Q & 1.004 & -0.4\% & 15.4\% & 35.2\% \\
nonlinear\_one\_hop & GHAR\_Q\_IV & GNNHAR1L\_Q\_IV & 1.001 & -0.1\% & 7.7\% & 40.7\% \\
linear\_multihop & GHAR2H\_M & GHAR3H\_M & 1.003 & -0.3\% & 33.0\% & 49.5\% \\
linear\_multihop & GHAR2H\_M\_IV & GHAR3H\_M\_IV & 1.001 & -0.1\% & 18.7\% & 39.6\% \\
linear\_multihop & GHAR\_M & GHAR2H\_M & 1.006 & -0.6\% & 20.9\% & 36.3\% \\
linear\_multihop & GHAR\_M\_IV & GHAR2H\_M\_IV & 0.999 & 0.1\% & 16.5\% & 51.6\% \\
gnn\_depth & GNNHAR1L\_M & GNNHAR2L\_M & 1.003 & -0.3\% & 17.6\% & 45.1\% \\
gnn\_depth & GNNHAR1L\_M\_IV & GNNHAR2L\_M\_IV & 1.026 & -2.6\% & 18.7\% & 42.9\% \\
gnn\_depth & GNNHAR1L\_Q & GNNHAR2L\_Q & 0.999 & 0.1\% & 3.3\% & 54.9\% \\
gnn\_depth & GNNHAR1L\_Q\_IV & GNNHAR2L\_Q\_IV & 0.999 & 0.1\% & 14.3\% & 62.6\% \\
gnn\_depth & GNNHAR2L\_M & GNNHAR3L\_M & 1.000 & -0.0\% & 11.0\% & 58.2\% \\
gnn\_depth & GNNHAR2L\_M\_IV & GNNHAR3L\_M\_IV & 0.978 & 2.2\% & 20.9\% & 54.9\% \\
gnn\_depth & GNNHAR2L\_Q & GNNHAR3L\_Q & 0.998 & 0.2\% & 11.0\% & 54.9\% \\
gnn\_depth & GNNHAR2L\_Q\_IV & GNNHAR3L\_Q\_IV & 1.007 & -0.7\% & 15.4\% & 37.4\% \\
gnn\_depth & GNNHAR3L\_M & GNNHAR4L\_M & 1.004 & -0.4\% & 7.7\% & 31.9\% \\
gnn\_depth & GNNHAR3L\_M\_IV & GNNHAR4L\_M\_IV & 1.006 & -0.6\% & 15.4\% & 53.8\% \\
gnn\_depth & GNNHAR3L\_Q & GNNHAR4L\_Q & 1.001 & -0.1\% & 9.9\% & 46.2\% \\
gnn\_depth & GNNHAR3L\_Q\_IV & GNNHAR4L\_Q\_IV & 0.996 & 0.4\% & 12.1\% & 61.5\% \\
gnn\_depth & GNNHAR4L\_M & GNNHAR5L\_M & 1.007 & -0.7\% & 4.4\% & 47.3\% \\
gnn\_depth & GNNHAR4L\_M\_IV & GNNHAR5L\_M\_IV & 0.998 & 0.2\% & 11.0\% & 60.4\% \\
gnn\_depth & GNNHAR4L\_Q & GNNHAR5L\_Q & 0.998 & 0.2\% & 13.2\% & 60.4\% \\
gnn\_depth & GNNHAR4L\_Q\_IV & GNNHAR5L\_Q\_IV & 0.998 & 0.2\% & 8.8\% & 57.1\% \\
\bottomrule
\end{tabular}
}

\vspace{2pt}\begin{minipage}{0.98\textwidth}\footnotesize A candidate/base ratio below one means the candidate has lower average QLIKE than the base.  Positive DM share is the share of tickers where the DM statistic favors the candidate.\end{minipage}
\end{table}

\begin{table}[H]
\centering
\caption{QLIKE DM comparison summary for SP500}
\label{tab:sp500-dm}
\tiny
\resizebox{\textwidth}{!}{%
\begin{tabular}{lrrrrrr}
\toprule
Type & Base & Candidate & Cand./base & Gain & p<0.05 share & Positive share \\
\midrule
graph\_linear & HAR\_M & GHAR\_M & 0.998 & 0.2\% & 17.6\% & 51.2\% \\
graph\_linear & HAR\_M\_IV & GHAR\_M\_IV & 0.999 & 0.1\% & 26.9\% & 56.1\% \\
graph\_linear & HAR\_Q & GHAR\_Q & 0.980 & 2.0\% & 70.6\% & 61.7\% \\
graph\_linear & HAR\_Q\_IV & GHAR\_Q\_IV & 1.008 & -0.8\% & 65.0\% & 40.3\% \\
nonlinear\_one\_hop & GHAR\_M & GNNHAR1L\_M & 10.943 & -994.3\% & 88.2\% & 4.2\% \\
nonlinear\_one\_hop & GHAR\_M\_IV & GNNHAR1L\_M\_IV & 11.211 & -1021.1\% & 92.2\% & 1.8\% \\
nonlinear\_one\_hop & GHAR\_Q & GNNHAR1L\_Q & 2.756 & -175.6\% & 86.4\% & 25.2\% \\
nonlinear\_one\_hop & GHAR\_Q\_IV & GNNHAR1L\_Q\_IV & 2.899 & -189.9\% & 87.8\% & 20.9\% \\
linear\_multihop & GHAR2H\_M & GHAR3H\_M & 1.000 & -0.0\% & 11.8\% & 49.2\% \\
linear\_multihop & GHAR2H\_M\_IV & GHAR3H\_M\_IV & 1.030 & -3.0\% & 58.1\% & 8.0\% \\
linear\_multihop & GHAR2H\_Q & GHAR3H\_Q & 0.999 & 0.1\% & 45.0\% & 60.4\% \\
linear\_multihop & GHAR2H\_Q\_IV & GHAR3H\_Q\_IV & 0.999 & 0.1\% & 73.3\% & 52.1\% \\
linear\_multihop & GHAR\_M & GHAR2H\_M & 1.000 & -0.0\% & 13.4\% & 48.6\% \\
linear\_multihop & GHAR\_M\_IV & GHAR2H\_M\_IV & 1.018 & -1.8\% & 25.2\% & 28.3\% \\
linear\_multihop & GHAR\_Q & GHAR2H\_Q & 1.017 & -1.7\% & 47.0\% & 33.2\% \\
linear\_multihop & GHAR\_Q\_IV & GHAR2H\_Q\_IV & 0.972 & 2.8\% & 61.2\% & 75.7\% \\
gnn\_depth & GNNHAR1L\_M & GNNHAR2L\_M & 1.031 & -3.1\% & 39.4\% & 36.7\% \\
gnn\_depth & GNNHAR1L\_M\_IV & GNNHAR2L\_M\_IV & 1.004 & -0.4\% & 29.2\% & 38.8\% \\
gnn\_depth & GNNHAR1L\_Q & GNNHAR2L\_Q & 1.038 & -3.8\% & 63.9\% & 23.6\% \\
gnn\_depth & GNNHAR1L\_Q\_IV & GNNHAR2L\_Q\_IV & 0.997 & 0.3\% & 58.8\% & 49.0\% \\
gnn\_depth & GNNHAR2L\_M & GNNHAR3L\_M & 0.976 & 2.4\% & 39.4\% & 46.5\% \\
gnn\_depth & GNNHAR2L\_M\_IV & GNNHAR3L\_M\_IV & 0.998 & 0.2\% & 23.2\% & 44.3\% \\
gnn\_depth & GNNHAR2L\_Q & GNNHAR3L\_Q & 0.963 & 3.7\% & 63.7\% & 70.6\% \\
gnn\_depth & GNNHAR2L\_Q\_IV & GNNHAR3L\_Q\_IV & 1.003 & -0.3\% & 52.3\% & 39.0\% \\
gnn\_depth & GNNHAR3L\_M & GNNHAR4L\_M & 1.010 & -1.0\% & 36.5\% & 45.0\% \\
gnn\_depth & GNNHAR3L\_M\_IV & GNNHAR4L\_M\_IV & 1.010 & -1.0\% & 34.7\% & 35.6\% \\
gnn\_depth & GNNHAR3L\_Q & GNNHAR4L\_Q & 1.269 & -26.9\% & 84.6\% & 7.6\% \\
gnn\_depth & GNNHAR3L\_Q\_IV & GNNHAR4L\_Q\_IV & 0.998 & 0.2\% & 42.5\% & 56.8\% \\
gnn\_depth & GNNHAR4L\_M & GNNHAR5L\_M & 0.984 & 1.6\% & 40.1\% & 51.0\% \\
gnn\_depth & GNNHAR4L\_M\_IV & GNNHAR5L\_M\_IV & 0.991 & 0.9\% & 33.4\% & 50.3\% \\
gnn\_depth & GNNHAR4L\_Q & GNNHAR5L\_Q & 0.788 & 21.2\% & 84.2\% & 88.4\% \\
gnn\_depth & GNNHAR4L\_Q\_IV & GNNHAR5L\_Q\_IV & 1.002 & -0.2\% & 48.6\% & 34.3\% \\
\bottomrule
\end{tabular}
}

\vspace{2pt}\begin{minipage}{0.98\textwidth}\footnotesize A candidate/base ratio below one means the candidate has lower average QLIKE than the base.  Positive DM share is the share of tickers where the DM statistic favors the candidate.\end{minipage}
\end{table}

\section{Regime, FVU, and Graph Diagnostics}

\begin{table}[H]
\centering
\caption{Best regime-specific loss ratios}
\label{tab:regime-best}
\scriptsize
\resizebox{\textwidth}{!}{%
\begin{tabular}{lrrrrrr}
\toprule
Universe & Regime & Metric & Best model & Ratio & Gain & Dates \\
\midrule
dow30 & calm\_bottom\_90pct & MSE & GHAR\_M\_IV & 0.973 & 2.7\% & 208 \\
dow30 & calm\_bottom\_90pct & QLIKE & GNNHAR2L\_Q\_IV & 0.947 & 5.3\% & 208 \\
dow30 & turbulent\_top\_10pct & MSE & GHAR\_M\_IV & 0.766 & 23.4\% & 24 \\
dow30 & turbulent\_top\_10pct & QLIKE & GHAR\_M\_IV & 0.664 & 33.6\% & 24 \\
sp100 & calm\_bottom\_90pct & MSE & GHAR3H\_M\_IV & 0.969 & 3.1\% & 210 \\
sp100 & calm\_bottom\_90pct & QLIKE & GNNHAR2L\_Q\_IV & 0.933 & 6.7\% & 210 \\
sp100 & turbulent\_top\_10pct & MSE & HAR\_M\_IV & 0.929 & 7.1\% & 24 \\
sp100 & turbulent\_top\_10pct & QLIKE & HAR\_Q\_IV & 0.921 & 7.9\% & 24 \\
sp500 & calm\_bottom\_90pct & MSE & HAR\_M\_IV & 0.968 & 3.2\% & 200 \\
sp500 & calm\_bottom\_90pct & QLIKE & GHAR\_M\_IV & 0.976 & 2.4\% & 200 \\
sp500 & turbulent\_top\_10pct & MSE & GHAR\_M\_IV & 0.944 & 5.6\% & 23 \\
sp500 & turbulent\_top\_10pct & QLIKE & GHAR\_M\_IV & 0.943 & 5.7\% & 23 \\
\bottomrule
\end{tabular}
}

\vspace{2pt}\begin{minipage}{0.98\textwidth}\footnotesize The current regime split uses the cross-sectional mean RV proxy because SPY is not present in these saved truth arrays or raw panels.  Zhang uses SPY RV, so this is method-aligned but not data-identical.\end{minipage}
\end{table}

\begin{table}[H]
\centering
\caption{Lowest FVU relative to HAR\_M by universe and regime}
\label{tab:fvu-best}
\scriptsize
\resizebox{\textwidth}{!}{%
\begin{tabular}{lrrrrrr}
\toprule
Universe & Regime & Model & Mean FVU & Median FVU & IV & Est. \\
\midrule
dow30 & calm\_bottom\_90pct & HAR\_M & 0.0000 & 0.0000 & no & MSE \\
dow30 & turbulent\_top\_10pct & HAR\_M & 0.0000 & 0.0000 & no & MSE \\
sp100 & calm\_bottom\_90pct & HAR\_M & 0.0000 & 0.0000 & no & MSE \\
sp100 & turbulent\_top\_10pct & HAR\_M & 0.0000 & 0.0000 & no & MSE \\
sp500 & calm\_bottom\_90pct & HAR\_M & 0.0000 & 0.0000 & no & MSE \\
sp500 & turbulent\_top\_10pct & HAR\_M & 0.0000 & 0.0000 & no & MSE \\
\bottomrule
\end{tabular}
}

\vspace{2pt}\begin{minipage}{0.98\textwidth}\footnotesize FVU is an effect-size diagnostic for how far a model's forecast function moves away from \(HAR_{\mathrm{M}}\); it is not a significance test.\end{minipage}
\end{table}

\begin{table}[H]
\centering
\caption{Available graph rollup}
\label{tab:graph-rollup}
\scriptsize
\resizebox{\textwidth}{!}{%
\begin{tabular}{lrrrrrrr}
\toprule
Universe & Loss group & Graphs & Nodes & Edges & Density & Mean degree & Diameter \\
\midrule
sp500 & MSE & 11 & 449.0 & 2238.8 & 0.022 & 9.97 & 10.73 \\
sp500 & QLIKE & 11 & 449.0 & 2238.8 & 0.022 & 9.97 & 10.73 \\
\bottomrule
\end{tabular}
}

\vspace{2pt}\begin{minipage}{0.98\textwidth}\footnotesize Only the S\&P 500 AutoDL package currently contains graph matrices in the paper-ready tree; Dow 30 and S\&P 100 forecasts remain analyzable, but graph structural summaries require saved graph matrices.\end{minipage}
\end{table}

\section{Comparison with Zhang et al.}

Zhang et al. (2024a) find that GHAR improves on HAR, GNNHAR1L or GNNHAR2L often improves further,
QLIKE training is useful especially under turbulent regimes, and very deep GNNHAR can deteriorate
because of over-smoothing.  Our Dow 30 result is directionally close once IV is included:
\texttt{GNNHAR2L\_Q\_IV} is the best QLIKE model and improves roughly \(11.1\%\) over \(HAR_{\mathrm{M}}\).
The S\&P 100 run is weaker for graph nonlinearity: IV helps, but the best QLIKE model is
\(HAR_Q\_IV\), and several IV graph/neural variants are statistically close by MCS.
The S\&P 500 AutoDL run is different again.  It is internally consistent over 449 selected
tickers, but the best QLIKE model is \texttt{GHAR\_M\_IV}, and QLIKE-trained neural models perform
poorly in the saved run.  This means the large-universe extension does not currently show a
stronger GNN improvement just because the graph has more nodes.  It instead suggests that
large-universe GNN training is more sensitive to optimization, scaling, and possibly over-smoothing.


\section{Current Scope Gaps}

The current local outputs support the one-day Zhang-style analysis.  The following pieces
cannot be reconstructed exactly from the saved forecasts alone:
\begin{itemize}
\item one-week and one-month horizon tables, because those require horizon-specific targets and reruns;
\item Zhang's exact SPY regime split, unless a SPY RV series is added and aligned;
\item exact hidden-state MAD from Zhang's Figure 7, because hidden representations were not saved;
\item the smaller-validation robustness table, because that changes the rolling train-validation split.
\end{itemize}
The generated CSV layer records these limitations explicitly in
\texttt{outputs/paper\_ready\_20260617/zhang\_style\_statistics/zhang\_statistics\_reference\_coverage.csv}.


\section{Generated Files}

The complete machine-readable results are stored under
\texttt{outputs/paper\_ready\_20260617/zhang\_style\_statistics/} and under each universe's
\texttt{zhang\_style\_statistics/} folder.  The most important cross-universe files are:
\begin{itemize}
\item \texttt{cross\_universe\_loss\_ratio\_summary.csv}
\item \texttt{cross\_universe\_mcs\_rollup.csv}
\item \texttt{cross\_universe\_pairwise\_qlike\_dm\_summary.csv}
\item \texttt{cross\_universe\_fvu\_by\_regime\_vs\_HAR\_M.csv}
\item \texttt{cross\_universe\_ticker\_alignment\_audit.csv}
\item \texttt{cross\_universe\_graph\_rollup.csv}
\end{itemize}

\section*{References}
\small
Andersen, T. G., Bollerslev, T., Diebold, F. X., and Ebens, H. (2001). The distribution of realized stock return volatility. \textit{Journal of Financial Economics}.\\
Barndorff-Nielsen, O. E., and Shephard, N. (2002). Econometric analysis of realized volatility and its use in estimating stochastic volatility models. \textit{Journal of the Royal Statistical Society: Series B}.\\
Bollerslev, T., Hood, B., Huss, J., and Pedersen, L. H. (2018). Risk everywhere: Modeling and managing volatility. \textit{Review of Financial Studies}. DOI: 10.1093/rfs/hhy041.\\
Busch, T., Christensen, B. J., and Nielsen, M. O. (2010). The role of implied volatility in forecasting future realized volatility and jumps in foreign exchange, stock, and bond markets. \textit{Journal of Econometrics}. DOI: 10.1016/j.jeconom.2010.03.014.\\
Chen, D., Lin, Y., Li, W., Li, P., Zhou, J., and Sun, X. (2020). Measuring and relieving the over-smoothing problem for graph neural networks from the topological view. \textit{Proceedings of the AAAI Conference on Artificial Intelligence}. DOI: 10.1609/aaai.v34i04.5747.\\
Corsi, F. (2009). A simple approximate long-memory model of realized volatility. \textit{Journal of Financial Econometrics}. DOI: 10.1093/jjfinec/nbp001.\\
Diebold, F. X., and Mariano, R. S. (1995). Comparing predictive accuracy. \textit{Journal of Business \& Economic Statistics}. DOI: 10.1080/07350015.1995.10524599.\\
Friedman, J., Hastie, T., and Tibshirani, R. (2008). Sparse inverse covariance estimation with the graphical lasso. \textit{Biostatistics}. DOI: 10.1093/biostatistics/kxm045.\\
Hansen, P. R., Lunde, A., and Nason, J. M. (2011). The model confidence set. \textit{Econometrica}. DOI: 10.3982/ECTA5771.\\
Kipf, T. N., and Welling, M. (2017). Semi-supervised classification with graph convolutional networks. \textit{International Conference on Learning Representations}. arXiv: 1609.02907.\\
Li, S. Z., and Tang, Y. (2021). Forecasting realized volatility: An automatic system using many features and many machine learning algorithms. \textit{SSRN Electronic Journal}. DOI: 10.2139/ssrn.3776915.\\
Lin, W. (2013). Agnostic notes on regression adjustments to experimental data: Reexamining Freedman's critique. \textit{The Annals of Applied Statistics}. DOI: 10.1214/12-AOAS583.\\
Patton, A. J. (2011). Volatility forecast comparison using imperfect volatility proxies. \textit{Journal of Econometrics}. DOI: 10.1016/j.jeconom.2010.03.034.\\
Zhang, C., Pu, X., Cucuringu, M., and Dong, X. (2024a). Forecasting realized volatility with spillover effects: Perspectives from graph neural networks. \textit{International Journal of Forecasting}. DOI: 10.1016/j.ijforecast.2024.09.002.\\
Zhang, C., Pu, X., Cucuringu, M., and Dong, X. (2024b). Graph-based methods for forecasting realized covariances. \textit{SSRN Electronic Journal}. DOI: 10.2139/ssrn.4274989.

\end{document}
